%% toc: P2 20 20.1,20.2,20.3,20.4
\section{Completeness and compactness without the Axiom of Choice}
\label{sec:p02-w02c-p02-023-choice-free}

The completeness theorem \cref{thm:fol-completeness}, the sequent-calculus
completeness theorem \cref{thm:fol-lk-completeness}, the compactness theorem
\cref{thm:fol-compactness} and the corollary that follows it on finite
structures each assume the Axiom of Choice in the ambient set theory, for an
arbitrary signature.
The hypothesis enters through the Henkin--Lindenbaum construction
(\cref{lem:henkin-lindenbaum}), which well-orders the formulas of the
witness-extended signature.  This section records which of the four
conclusions can be obtained with no choice principle at all, proves the
compactness cases that are not proved elsewhere in the handbook, and shows
that the restriction to countable signatures cannot simply be removed.

\emph{Inputs.}  A first-order signature \(\Lang\) in the sense of
\cref{sec:fol-syntax-semantics-01}, a theory \(T\subseteq\operatorname{Sent}(\Lang)\) and
a sentence \(\sigma\), with the semantics of \cref{ch:fol-theories} and the
natural-deduction calculus \(\mathsf{ND}_{\mathsf{FOL}}\).  \emph{Conclusions.}  For
countable \(\Lang\): both forms of compactness and the non-axiomatizability
of finiteness, in \(\mathsf{ZF}\).  For arbitrary \(\Lang\): a reduction
showing that compactness for every signature yields the Boolean Prime Ideal
Theorem.

\subsection{Assumptions and the choice-free counterparts}
\label{sec:p02-w02c-p02-023-assumptions}

\begin{hbconvention}{def:p02-w02c-p02-023-zf}{Ambient theory and countable signatures}
In this section the ambient set theory is \(\mathsf{ZF}\), without the Axiom
of Choice, unless a statement says otherwise.  A signature
\(\Lang\) with function symbols \(\mathcal F\) and relation symbols \(\mathcal R\)
is \emph{countable} when the set \(\mathcal F\cup\mathcal R\) of its function
and relation symbols is countable in the sense of
\cref{sec:set-prerequisites-05}: some injection
\(\mathcal F\cup\mathcal R\to\N\) exists.  Finite signatures, the empty one
included, are countable.  Countability is a property witnessed by an
injection; listing the symbols in a display does not supply one for an
infinite symbol set.  A \emph{model} of a theory, satisfiability and the
consequence relation \(\entails\) are those of
\cref{def:fol-theory-satisfiability} and the semantic-consequence definition
of \cref{sec:fol-semantics-02}; \(T\derives_{\mathsf{ND}_{\mathsf{FOL}}}\varphi\) means that
some \(\mathsf{ND}_{\mathsf{FOL}}\)-derivation of \(\varphi\) has all its open assumptions in
\(T\).
\end{hbconvention}

The choice-free counterparts of the four inherited results are as follows.
Each is a separate result with its own hypothesis; none of them changes the
statement or the hypothesis of the inherited result it accompanies.
\begin{enumerate}
  \item \emph{Natural-deduction completeness.}  For countable \(\Lang\),
        \cref{cor:fol-completeness-countable} proves
        \(\Gamma\entails\varphi\Rightarrow\Gamma\derives_{\mathsf{ND}_{\mathsf{FOL}}}\varphi\)
        with no use of the Axiom of Choice.
  \item \emph{Sequent-calculus completeness.}  For every signature and every
        finite sequent, \cref{thm:p07-eqlk-completeness} proves that validity
        is equivalent to derivability, with Cut, on the equality sequent
        calculus \texttt{sheet:lk-eq}, which is the calculus \(\mathsf{LK}_{\mathsf{FOL}}\)
        of \cref{sec:fol-deduction-02} treated in
        \cref{thm:fol-lk-completeness}, without the Axiom of Choice.  Its
        proof reduces the sequent to the finite signature of the symbols
        occurring in it and applies \cref{cor:fol-completeness-countable}.
        For the equality-free calculus the same reduction is
        \cref{thm:p02-w02b-p02-019-choice-free}.
  \item \emph{Compactness.}  For countable \(\Lang\),
        \cref{thm:p02-w02c-p02-023-countable-compactness} below.
  \item \emph{Finite structures.}
        \Cref{thm:p02-w02c-p02-023-finite-structures} below.
\end{enumerate}

\subsection{Compactness for countable signatures}
\label{sec:p02-w02c-p02-023-countable}

\begin{hbtheorem}{thm:p02-w02c-p02-023-countable-compactness}{proved=annex:p02-w02c-p02-023-compact}{Countable compactness without choice}
Let \(\Lang\) be a countable signature.  For every
\(T\subseteq\operatorname{Sent}(\Lang)\) and every \(\sigma\in\operatorname{Sent}(\Lang)\):
\begin{enumerate}
  \item \(T\) has a model iff every finite subset of \(T\) has a model; and
  \item if \(T\entails\sigma\), then some finite \(T_0\subseteq T\)
        satisfies \(T_0\entails\sigma\).
\end{enumerate}
No use is made of the Axiom of Choice.
\end{hbtheorem}

The proof does not build a model directly.  It shows that \(T\) is
consistent, by soundness, and then uses the choice-free completeness of
\cref{cor:fol-completeness-countable} with \(\varphi=\bottom\) in
contrapositive form: an unsatisfiable \(T\) would derive \(\bottom\).  The
model exists because the Henkin construction of that corollary runs along
one explicit enumeration of the formulas of the witness-extended signature.

\begin{hbcorollary}{thm:p02-w02c-p02-023-finite-structures}{proved=annex:p02-w02c-p02-023-finite}{Finiteness unaxiomatizable without choice}
No use is made of the Axiom of Choice.  In the empty signature, no theory
\(T\) has as its models exactly the structures with finite domain.
\end{hbcorollary}

\begin{hbmethod}{def:p02-w02c-p02-023-method}{Deciding which compactness theorem an argument may cite}
\begin{enumerate}
  \item Name the signature \(\Lang\) in which the theory is written,
        including every constant symbol added for the argument.
  \item Exhibit an injection of its symbol set into \(\N\).  For a finite
        signature, or one obtained from a finite signature by adding the
        constants \(c_n\) for \(n\in\N\), the injection can be written
        down.
  \item If such an injection is exhibited, cite
        \cref{thm:p02-w02c-p02-023-countable-compactness}; the argument then
        uses no choice principle at this step.
  \item Otherwise cite \cref{thm:fol-compactness}, whose hypothesis
        includes the Axiom of Choice, and record that hypothesis among the
        assumptions of the conclusion.
  \item Check the remaining steps separately.  A choice-free compactness
        step does not make a later appeal to the L\"owenheim--Skolem
        theorems, \cref{thm:downward-ls,thm:upward-ls}, choice-free: both
        assume the Axiom of Choice.
\end{enumerate}
\end{hbmethod}

\begin{hbexample}{def:p02-w02c-p02-023-example}{A nonstandard model of true arithmetic in ZF}
Let \(\Lang_A\), \(\mathcal N_{\mathrm{ar}}\), the new constant \(c\) and
\[
 T_1=\operatorname{Th}(\mathcal N_{\mathrm{ar}})\cup
 \{\neg\,c=\bar n:n\in\N\}
\]
be as in \cref{thm:p02-w02-p02-033-nonstandard-direct}.  The signature
\(\Lang_A\cup\{c\}\) is finite, so an injection into \(\N\) is given by
numbering its symbols, and \cref{thm:p02-w02c-p02-023-countable-compactness}
applies.

\emph{Finite subsets.}  Let \(\Delta\subseteq T_1\) be finite.  Put
\(m_\Delta=1+\max\{n:\neg\,c=\bar n\in\Delta\}\), and \(m_\Delta=0\) if no
such sentence lies in \(\Delta\).  Clause~1 of
\cref{thm:p02-w02-p02-033-nonstandard-direct} shows that the expansion of
\(\mathcal N_{\mathrm{ar}}\) interpreting \(c\) as \(m_\Delta\) satisfies
\(\Delta\).  For instance, for
\(\Delta=\{\neg\,c=\bar 0,\neg\,c=\bar 3\}\) we get \(m_\Delta=4\), and
\(4\) differs from \(0\) and from \(3\).

\emph{Conclusion.}  Clause~1 of
\cref{thm:p02-w02c-p02-023-countable-compactness} gives a model
\(\mathcal K\) of \(T_1\), with no use of the Axiom of Choice.  For each
\(n\), \(\mathcal K\entails\neg\,c=\bar n\), so the element
\(c^{\mathcal K}\) is distinct from the value in \(\mathcal K\) of every
numeral, while the \(\Lang_A\)-reduct of \(\mathcal K\) satisfies every
sentence true in \(\mathcal N_{\mathrm{ar}}\) by
\cref{lem:reduct-invariance}.

The example obtains the existence of such a model only.  Clause~2 of
\cref{thm:p02-w02-p02-033-nonstandard-direct}, a model with countable
domain, is reached there through \cref{thm:downward-ls}, which assumes the
Axiom of Choice; the compactness step above does not remove that
assumption.
\end{hbexample}

\subsection{The boundary at arbitrary signatures}
\label{sec:p02-w02c-p02-023-boundary}

The proof of \cref{thm:p02-w02c-p02-023-countable-compactness} needs an
enumeration of the formulas of the witness-extended signature, and a
countable signature supplies one explicitly.  For an arbitrary signature no
such enumeration need exist without some choice principle.  The next result
measures what an arbitrary-signature version of clause~1 would have to
supply.

\begin{hbtheorem}{thm:p02-w02c-p02-023-bpi}{proved=annex:p02-w02c-p02-023-bpi}{Compactness for all signatures gives BPI}
Work in \(\mathsf{ZF}\).  Suppose that clause~1 of \cref{thm:fol-compactness}
holds for every signature, with no choice hypothesis: for every signature
\(\Lang\) and every \(T\subseteq\operatorname{Sent}(\Lang)\), if every finite subset of
\(T\) has a model then \(T\) has a model.  Then the propositional
compactness principle of \cref{def:p01-bpi-propositional-compactness} holds
for every set \(V\) of variables, and hence, by
\cref{thm:p01-bpi-propositional-compactness}, so does each of the
equivalent forms of the Boolean Prime Ideal Theorem in
\cref{thm:p01-bpi-equivalence}.
\end{hbtheorem}

The reduction uses one relation symbol \(R_v\) of arity zero for each
propositional variable \(p_v\).  A structure for these symbols determines a
valuation and conversely, and the translation of a propositional formula is a
sentence true in a structure exactly when the formula is true under the
corresponding valuation.

\begin{hbwarning}{def:p02-w02c-p02-023-limits}{What the choice-free results do not give}
\begin{enumerate}
  \item \Cref{thm:p02-w02c-p02-023-countable-compactness} concerns countable
        signatures only.  It does not show that the Axiom of Choice is
        needed for uncountable ones;
        \cref{thm:p02-w02c-p02-023-bpi} shows only that a choice-free proof
        for every signature would also prove the Boolean Prime Ideal Theorem
        in \(\mathsf{ZF}\).  The relation between that theorem and the full
        Axiom of Choice is the subject of
        \cref{thm:p01-bpi-choice-separation}.
  \item The countable theorem gives a model but no bound on its size and no
        description of it beyond the sentences it satisfies.  Size bounds
        come from \cref{thm:downward-ls}, which assumes the Axiom of Choice.
  \item \Cref{thm:p07-eqlk-completeness} is a statement about finite
        sequents.  It does not give completeness for an infinite premise set
        over an uncountable signature; that case is
        \cref{thm:fol-completeness}, with its Choice hypothesis.
  \item The statements of \cref{thm:fol-completeness},
        \cref{thm:fol-lk-completeness}, \cref{thm:fol-compactness} and the
        corollary on finite structures keep their Choice hypothesis.  An
        argument that cites one of them carries that hypothesis even when
        its signature is countable, unless it cites the choice-free result
        instead.
\end{enumerate}
\end{hbwarning}

\subsection{Relationships and sources}
\label{sec:p02-w02c-p02-023-sources}

The four inherited results and their proofs are those of the frozen
reference \cite[lines 4713--4967, sections ``Completeness'' and
``Compactness'']{FrozenFormalLogicCompactnessApplications}, where the
theorems at lines 4716--4726, 4832--4843 and 4924--4936 and the corollary at
lines 4954--4958 each state the Axiom of Choice as a hypothesis, and the
choice-free countable completeness corollary is stated and proved at lines
4760--4830.  \Cref{thm:p02-w02c-p02-023-countable-compactness},
\cref{thm:p02-w02c-p02-023-finite-structures} and
\cref{thm:p02-w02c-p02-023-bpi}, with their proofs, are editorial results
built from the definitions and results of that reference and of
\cref{sec:p01-bpi-compactness-bridge}; they are not attributed to an
external presentation.  The proof of
\cref{thm:p02-w02c-p02-023-countable-compactness} follows the proof of
clause~1 of \cref{thm:fol-compactness} up to the consistency of \(T\) and then
replaces the Henkin construction under Choice by
\cref{cor:fol-completeness-countable}.  The choice-dependency table of
\cref{sec:p06-cardinals-choice} lists the principle each related theorem
uses.
